\chapter{1977年山西卷}

\begin{problem}
化简下列各题.

\begin{enumerate}
\item \( \frac{1}{\sqrt{2}+1}-\frac{1}{\sqrt{2}-1} \)；

\item \( \frac{(a^2-b^2)^3}{a^3+b^3}\div\frac{(b+a)^2}{a^2-ab+b^2}\times\frac{1}{(b-a)^3} \)；

\item \( \frac{2\lg 6-\lg 3}{1+\frac{1}{2}\lg 0.36+\frac{1}{3}\lg 8} \)；

\item \( \frac{x^4+x^2+1}{x^3-1}+\frac{x(x-2)-2(x-2)}{x^2-3x+2} \).
\end{enumerate}
\end{problem}

\begin{solution}
\begin{enumerate}
\item 有理化分母，得
\[
\frac{1}{\sqrt{2}+1}-\frac{1}{\sqrt{2}-1}
=\frac{\left(\sqrt{2}-1\right)-\left(\sqrt{2}+1\right)}{\left(\sqrt{2}+1\right)\left(\sqrt{2}-1\right)}
=\frac{-2}{2-1}=-2
\]

\item 在实数范围内，原式有意义需要 \(a^3+b^3\ne0\)、\(a^2-ab+b^2\ne0\)、\(\dfrac{(a+b)^2}{a^2-ab+b^2}\ne0\) 且 \(a\ne b\). 又因为 \(a^2-ab+b^2=\dfrac{1}{2}\left((a-b)^2+a^2+b^2\right)\)，所以在这些条件下等价于 \(a+b\ne0\) 且 \(a\ne b\). 利用 \(a^2-b^2=(a+b)(a-b)\)、\(a^3+b^3=(a+b)(a^2-ab+b^2)\) 以及 \((b-a)^3=-(a-b)^3\)，得
\[
\frac{(a^2-b^2)^3}{a^3+b^3}\cdot\frac{a^2-ab+b^2}{(a+b)^2}\cdot\frac{1}{(b-a)^3}
=\frac{(a+b)^3(a-b)^3}{(a+b)(a^2-ab+b^2)}\cdot\frac{a^2-ab+b^2}{(a+b)^2}\cdot\frac{1}{-(a-b)^3}=-1
\]

\item 因为 \(0.36=0.6^2\)、\(8=2^3\) 且 \(1=\lg 10\)，所以分子为 \(2\lg6-\lg3=\lg\dfrac{6^2}{3}=\lg12\)，分母为
\[
1+\frac{1}{2}\lg0.36+\frac{1}{3}\lg8
=\lg10+\lg0.6+\lg2
=\lg\left(10\times0.6\times2\right)=\lg12
\]
因此原式 \(=\dfrac{\lg12}{\lg12}=1\).

\item 原式有意义时，\(x^3-1\ne0\) 且 \(x^2-3x+2\ne0\)，即 \(x\ne1\)，\(2\). 在此条件下，
\[
\begin{gathered}
\text{原式}=\frac{(x^2+1)^2-x^2}{(x-1)(x^2+x+1)}+\frac{(x-2)^2}{(x-2)(x-1)}\\
=\frac{(x^2-x+1)(x^2+x+1)}{(x-1)(x^2+x+1)}+\frac{x-2}{x-1}
=\frac{x^2-x+1+x-2}{x-1}\\
=\frac{x^2-1}{x-1}=x+1.
\end{gathered}
\]
\end{enumerate}
\end{solution}

\begin{problem}
作下列各题.

\begin{enumerate}
\item 当 \(a\)，\(b\)，\(c>0\) 时，求证：\( \log_a b\cdot\log_b c\cdot\log_c a=1 \).

\item 已知 \( y=\lg \frac{1}{1-\sqrt{1-x}} \)，求定义域.

\item 汽车往返甲乙两地之间，上行速度为 \(30\text{ 公里/小时}\)，下行速度为 \(60\text{ 公里/小时}\). 求往返平均速度.
\end{enumerate}
\end{problem}

\begin{solution}
\begin{enumerate}
\item 题中三个对数有定义时，除 \(a\)，\(b\)，\(c>0\) 外还需 \(a\)，\(b\)，\(c\ne1\). 在此条件下，换底公式给出
\(\log_a b=\dfrac{\lg b}{\lg a}\)，\(\log_b c=\dfrac{\lg c}{\lg b}\)，\(\log_c a=\dfrac{\lg a}{\lg c}\). 因而
\[
\log_a b\cdot\log_b c\cdot\log_c a
=\frac{\lg b}{\lg a}\cdot\frac{\lg c}{\lg b}\cdot\frac{\lg a}{\lg c}=1
\]

\item 根式有意义且对数的真数为正，必须满足
\[
\begin{cases}
1-x\ge0,\\
\dfrac{1}{1-\sqrt{1-x}}>0.
\end{cases}
\]
由于分子为正，第二个条件等价于 \(1-\sqrt{1-x}>0\). 于是 \(\sqrt{1-x}<1\)，从而 \(x>0\)，再与 \(x\le1\) 合并，得到定义域为 \(0<x\le1\).

\item 设甲、乙两地之间的距离为 \(m\text{ 公里}\)，其中 \(m>0\). 则上行和下行的路程都为 \(m\)，所用时间分别为 \(t_1=\dfrac{m}{30}\text{ 小时}\) 和 \(t_2=\dfrac{m}{60}\text{ 小时}\). 往返平均速度等于总路程除以总时间，所以
\[
\overline{V}=\frac{2m}{t_1+t_2}=\frac{2m}{\frac{m}{30}+\frac{m}{60}}=\frac{2m}{\frac{m}{20}}=40\text{ 公里/小时}
\]
答：往返平均速度为 \(40\text{ 公里/小时}\).
\end{enumerate}
\end{solution}

\begin{problem}
试证梯形的中位线平行于两底，且等于两底和的一半.

已知：如图，梯形 \(ABCD\)，\( AD\parallel BC \)，\(M\)，\(N\) 分别为 \(AB\)，\(DC\) 之中点.

求证：

\begin{enumerate}
\item \(MN\parallel AD\)，\(MN\parallel BC\)；
\item \(MN=\frac{1}{2}(AD+BC)\).
\end{enumerate}

\begin{center}
\bitmapfigure[float=inline,width=4.2cm]{img/1977/shanxi/q3_fig1.png}
\end{center}
\end{problem}

\begin{solution}
\begin{enumerate}
\item 连接 \(DM\)，并延长 \(DM\) 交底边 \(BC\) 的延长线于 \(E\). 因为 \(M\) 是 \(AB\) 的中点，所以 \(AM=BM\). 又因为直线 \(DM\) 与直线 \(AB\) 相交，故 \(\angle EMB=\angle AMD\)；由 \(AD\parallel BE\)，得 \(\angle MBE=\angle MAD\). 因而 \(\triangle EBM\) 与 \(\triangle DAM\) 有两角及其夹边分别相等，所以 \(\triangle EBM\cong\triangle DAM\). 从而 \(BE=AD\)，\(EM=DM\). 由于 \(D\)，\(M\)，\(E\) 三点共线，\(EM=DM\) 说明 \(M\) 是 \(DE\) 的中点；又因 \(N\) 是 \(DC\) 的中点，所以在 \(\triangle DEC\) 中，\(MN\) 是中位线. 因此 \(MN\parallel EC\). 由于 \(E\)，\(B\)，\(C\) 三点共线，\(EC\) 与 \(BC\) 为同一直线，所以 \(MN\parallel BC\)；再结合已知的 \(AD\parallel BC\)，得到 \(MN\parallel AD\).

\item 由三角形中位线定理，\(MN=\dfrac{1}{2}EC\). 按图中点的次序，\(E\)，\(B\)，\(C\) 三点共线且 \(B\) 在 \(EC\) 上，因此 \(EC=EB+BC\). 又由全等三角形得到 \(EB=AD\)，所以
\[
MN=\frac{1}{2}EC=\frac{1}{2}(EB+BC)=\frac{1}{2}(AD+BC)
\]
\end{enumerate}
\end{solution}

\begin{problem}
有一块形状为直角三角形的白铁皮，其一直角边和斜边分别为 \(6\mathrm{~dm}\) 和 \(10\mathrm{~dm}\). 若从这块白铁皮剪一块最大圆材料，求这圆材料的面积有多大？（\(\pi\) 取 \(3.14\)）
\end{problem}

\begin{solution}
设直角三角形的两直角边为 \(a\)，\(b\)，斜边为 \(c\). 由勾股定理，另一条直角边为 \(b=\sqrt{10^2-6^2}=8\mathrm{~dm}\)，所以三角形面积为 \(S=\dfrac{1}{2}\times6\times8=24\mathrm{~dm}^2\)，半周长为 \(s=\dfrac{6+8+10}{2}=12\mathrm{~dm}\).

设任意一个完全位于三角形内的圆半径为 \(\rho\)，圆心到三边的距离分别为 \(d_1\)，\(d_2\)，\(d_3\). 因为圆在三角形内，\(d_i\ge\rho\)（\(i=1\)，\(2\)，\(3\)）. 以圆心为顶点把三角形分成三个三角形，得
\[
S=\frac{1}{2}(6d_1+8d_2+10d_3)\ge\frac{1}{2}(6+8+10)\rho=12\rho
\]
所以 \(\rho\le\dfrac{S}{s}=2\mathrm{~dm}\). 三角形的内切圆半径满足 \(S=rs\)，故 \(r=\dfrac{24}{12}=2\mathrm{~dm}\)，此圆达到上述上界，因而它就是能剪出的最大圆材料.

因此圆材料的面积为
\[
\pi r^2\approx3.14\times2^2=12.56\mathrm{~dm}^2
\]
答：剪得的最大圆材料的面积约为 \(12.56\mathrm{~dm}^2\).
\end{solution}

\begin{problem}
证明恒等式 \( \frac{1+\sin(180^\circ-\alpha)-\cos(180^\circ+\alpha)+\sin2\alpha}{1-\sin(180^\circ+\alpha)+\cos(\alpha-360^\circ)}=\sin\alpha+\cos\alpha \).
\end{problem}

\begin{solution}
在等式左边有意义的条件下，利用诱导公式
\(\sin(180^\circ-\alpha)=\sin\alpha\)，\(\cos(180^\circ+\alpha)=-\cos\alpha\)，\(\sin(180^\circ+\alpha)=-\sin\alpha\)，\(\cos(\alpha-360^\circ)=\cos\alpha\)，左边化为
\[
\frac{1+\sin\alpha+\cos\alpha+\sin2\alpha}{1+\sin\alpha+\cos\alpha}
\]
令 \(u=\sin\alpha+\cos\alpha\)，则由 \(\sin^2\alpha+\cos^2\alpha=1\) 和二倍角公式，\(u^2=1+\sin2\alpha\). 因而
\[
\frac{1+\sin\alpha+\cos\alpha+\sin2\alpha}{1+\sin\alpha+\cos\alpha}
=\frac{u^2+u}{1+u}=\frac{u(1+u)}{1+u}=u=\sin\alpha+\cos\alpha
\]
这里分母有意义正好保证 \(1+u\ne0\)，所以约分合法，恒等式得证.
\end{solution}

\begin{problem}
\(\triangle ABC\) 的三个内角是 \(\alpha\)，\(\beta\)，\(\gamma\)，而 \(\tan\alpha\)，\(\tan\beta\) 是方程 \(x^2-5x+6=0\) 的两个根. 求 \(\gamma\).
\end{problem}

\begin{solution}
因式分解方程得 \(x^2-5x+6=(x-2)(x-3)=0\)，所以 \(\tan\alpha\) 和 \(\tan\beta\) 分别为 \(2\) 和 \(3\)，从而 \(\tan\alpha+\tan\beta=5\)，\(\tan\alpha\tan\beta=6\). 由于三角形内角在 \(0^\circ\) 与 \(180^\circ\) 之间且正切值为正，\(\alpha\)、\(\beta\) 都是锐角，故 \(\alpha+\beta=180^\circ-\gamma\). 于是
\[
\tan\gamma=-\tan(\alpha+\beta)=-\frac{\tan\alpha+\tan\beta}{1-\tan\alpha\tan\beta}=-\frac{5}{1-6}=1
\]
又 \(0^\circ<\gamma<180^\circ\)，且 \(\tan\gamma>0\)，所以 \(\gamma\) 为锐角，得到 \(\gamma=45^\circ\).
\end{solution}

\begin{problem}
抛物线的顶点在原点，对称轴为 \(y\) 轴，过点 \((1,-1)\)，求抛物线方程，焦点坐标，准线方程，开口方向.
\end{problem}

\begin{solution}
顶点在原点且对称轴为 \(y\) 轴的抛物线可设为 \(x^2=4ay\). 代入已知点 \((1,-1)\)，得 \(1=4a(-1)\)，所以 \(a=-\dfrac{1}{4}\). 因而抛物线方程为 \(x^2=-y\). 对标准式 \(x^2=4ay\)，焦点为 \((0,a)\)，准线为 \(y=-a\)，所以焦点坐标为 \(\left(0,-\dfrac{1}{4}\right)\)，准线方程为 \(y=\dfrac{1}{4}\). 因为 \(a<0\)，抛物线开口向下.
\end{solution}

\begin{problem}
当 \( k \) 取何值时，直线 \( kx-y-3=0 \) 与圆 \( 2x^2+2y^2-9=0 \) 相切，并求出切点坐标.
\end{problem}

\begin{solution}
圆的方程可化为 \(x^2+y^2=\dfrac{9}{2}\)，所以圆心为原点，半径为 \(r=\dfrac{3}{\sqrt{2}}\). 直线 \(kx-y-3=0\) 到原点的距离为
\(d=\frac{3}{\sqrt{k^2+1}}\).
直线与圆相切当且仅当 \(d=r\)，故
\(\frac{3}{\sqrt{k^2+1}}=\frac{3}{\sqrt2}\)，\(\Longrightarrow\)，\(k^2+1=2\)，\(\Longrightarrow\)，\(k=\pm1\).

切点是圆心到切线的垂足. 当 \(k=1\) 时，切线为 \(y=x-3\)，过原点的半径垂直于它，故所在直线为 \(y=-x\). 联立得 \(x=\dfrac{3}{2}\)，\(y=-\dfrac{3}{2}\)，切点为 \(\left(\dfrac{3}{2},-\dfrac{3}{2}\right)\). 当 \(k=-1\) 时，切线为 \(y=-x-3\)，过原点的半径所在直线为 \(y=x\). 联立得 \(x=-\dfrac{3}{2}\)，\(y=-\dfrac{3}{2}\)，切点为 \(\left(-\dfrac{3}{2},-\dfrac{3}{2}\right)\).
\end{solution}

\section{参考题}

\begin{problem}
已知 \( f(x)=x\ln(2-x)+\sqrt{1+3x^2} \)，求 \( f'(1) \).
\end{problem}

\begin{solution}
函数的定义域为 \(x<2\)，且 \(1\) 在定义域内. 由乘积法则、链式法则，
\[
f'(x)=\ln(2-x)+x\left(-\frac{1}{2-x}\right)+\frac{3x}{\sqrt{1+3x^2}}=\ln(2-x)-\frac{x}{2-x}+\frac{3x}{\sqrt{1+3x^2}}
\]
代入 \(x=1\)，得
\[
f'(1)=\ln1-\frac{1}{2-1}+\frac{3}{\sqrt{1+3}}=0-1+\frac{3}{2}=\frac{1}{2}
\]
\end{solution}

\begin{problem}
一水库的水闸为一等腰梯形，上底为 \(8\text{ 米}\)，下底为 \(4\text{ 米}\)，高为 \(10\text{ 米}\). 求当水面与上底相齐时水闸所受的压力.（已知水的比重为 \(1\text{ 吨重/立方米}\)）

\bitmapfigure[float=inline,width=5.8cm]{img/1977/shanxi/q15_fig1.png}
\end{problem}

\begin{solution}
取水面为深度 \(x=0\)，向下为正方向，令 \(x\) 的单位为米. 水的比重为 \(1\text{ 吨重/立方米}\)，所以深度为 \(x\text{ 米}\) 处的压强为 \(P=x\text{ 吨重/平方米}\). 在深度 \(x\) 到 \(x+\mathrm{d}x\) 之间取一条水平薄片，则其宽度随深度线性变化，由上底宽 \(8\) 米、下底宽 \(4\) 米和高 \(10\) 米可得
\(w(x)=8-\frac{8-4}{10}x=8-\frac{2}{5}x\text{ 米}\). 薄片面积为 \(w(x)\mathrm{d}x\)，所受压力为 \(\mathrm{d}F=P\,w(x)\mathrm{d}x=x\left(8-\frac{2}{5}x\right)\mathrm{d}x\).
把深度从 \(0\) 米累加到 \(10\) 米，得到水闸所受总压力
\[
F=\int_0^{10}x\left(8-\frac{2}{5}x\right)\mathrm{d}x
=\left(4x^2-\frac{2}{15}x^3\right)\Big|_0^{10}
=400-\frac{2000}{15}=\frac{800}{3}\text{ 吨重}
\]
答：水闸所受的压力为 \(\dfrac{800}{3}\text{ 吨重}\).
\end{solution}
